\%============================================================
\chapter{Bai Hoc: Ban Be Tri Ky --- Nguoi Hieu Ta La Vo Gia (True Friendship --- A Priceless Understanding)}
\begin{center}
  \textit{\color{truyengold}``Ban be la nguoi biet tat ca ve ban va van yeu quy ban.''}\\[4pt]
  \textit{(A friend is someone who knows all about you and still loves you.)}\\[4pt]
  \small\color{truyenblue}--- Co Kim Dong Tay ---
\end{center}
\ngancach

\section{Ba Ya Va Tu Ky --- Am Nhac Tri Am}
\begin{truyen}{Ba Ya Va Tu Ky --- Am Nhac Tri Am}{Truyen Co Trung Hoa --- TK 5 TCN}
\chuhoa{B}{}B a Ya la mot nghe sy dan tranh bac thay thoi Xuan Thu; ong co the bieu hien moi cam xuc qua tieng dan.\\
\textit{(Boya was a master zither player of the Spring and Autumn period; he could express every emotion through his music.)}

Nhung chi mot nguoi hieu duoc am nhac cua ong --- Tu Ky, mot tieu phu theo doi ong tu rung ve.\\
\textit{(But only one person truly understood his music --- Ziqi, a woodcutter who followed him back from the forest.)}

Khi Ba Ya dan ve nui cao, Tu Ky noi: 'Tiet tau nhu Nui Thai Son vut len troi.' Khi dan nguoc, Tu Ky noi: 'Nuoc cuon cua Truong Giang.'\\
\textit{(When Boya played of high mountains, Ziqi said: 'The rhythm soars like Mount Tai.' When he played rivers, Ziqi said: 'The flow of the Yangtze.')}

Ba Ya kinh ngac: giua hai nguoi khac biet xa hoi hoan toan, trai tim ngu am nhu mot.\\
\textit{(Boya was astonished: between two people of completely different social standing, their musical hearts beat as one.)}

Khi Tu Ky mat, Ba Ya dap vo dan tranh: 'Tri am da mat, con ai nua de ta dan cho?'\\
\textit{(When Ziqi died, Boya smashed his zither: 'My soul-friend is gone; who remains for whom I should play?')}

\end{truyen}
\begin{baihoc}
  Tri am --- nguoi hieu trai tim ban --- quy hon van vang; thieu tri am, tat ca tai nang tro nen co don va vo nghia.\\
  \textit{(A soul-friend --- someone who understands your heart --- is worth more than ten thousand gold; without one, all talent becomes lonely and meaningless.)}
\end{baihoc}

\section{Orville Va Wilbur Wright --- Anh Em Cung Uoc Mo}
\begin{truyen}{Orville Va Wilbur Wright --- Anh Em Cung Uoc Mo}{Wright Brothers --- Hoa Ky 1903}
\chuhoa{O}{}O rville va Wilbur Wright khong co bang dai hoc, khong co tai tro chinh phu, khong co ekip ky su.\\
\textit{(Orville and Wilbur Wright had no university degrees, no government funding, no team of engineers.)}

Ho chi co nhau --- va cung mot giac mo dien ro rang ngay la: nguoi co the bay.\\
\textit{(They only had each other --- and one dream that seemed insane at the time: that people could fly.)}

Khi Wilbur thu kheo, Orville ghi chep ky luong; khi Orville nan long, Wilbur la nguoi dong vien.\\
\textit{(When Wilbur experimented, Orville took detailed notes; when Orville lost heart, Wilbur was there to encourage.)}

Ngay 17/12/1903, may bay mot dong co dau tien cat canh duoc 12 giay tai Kitty Hawk, Bac Carolina.\\
\textit{(On December 17, 1903, the first powered aircraft flew for 12 seconds at Kitty Hawk, North Carolina.)}

Ho khong tranh nhau ai la nguoi bay truoc --- ho quay phiem khi nao; nguoi ban thang duoc bay truoc.\\
\textit{(They didn't argue over who would fly first --- they flipped a coin; the winner flew first.)}

\end{truyen}
\begin{baihoc}
  Ban be that su khong tranh gianh cong trang --- ho chia se no; va chinh su chia se do nhan doi suc manh cho ca hai.\\
  \textit{(True friends don't compete for credit --- they share it; and that sharing doubles the strength of both.)}
\end{baihoc}

\section{Marie Curie Va Nguoi Ban Gay Dung Su Nghiep}
\begin{truyen}{Marie Curie Va Nguoi Ban Gay Dung Su Nghiep}{Marie Curie --- Ba Lan/Phap 1891}
\chuhoa{K}{}K hi Marie Sklodowska roi Ba Lan den Paris hoc, ba gan nhu khong co tien --- chi co sach va mot ban la Bronya chi em.\\
\textit{(When Marie Sklodowska left Poland for Paris to study, she had almost no money --- only books and the friend Bronya, her own sister.)}

Theo ke hoach cua hai chi em: Bronya hoc y truoc, Marie lam viec gui tien ho tro; xong Bronya lan luot giup Marie.\\
\textit{(By arrangement between the two sisters: Bronya studied medicine first while Marie worked and sent money; then Bronya would help Marie in turn.)}

Ke hoach nay keo dai nam nam --- cu bo cuoc thi muon; nhung ca hai giu loi hua.\\
\textit{(The plan stretched five years --- either could have abandoned it anytime; but both kept their promise.)}

Marie tro thanh nguoi phu nu dau tien nhan giai Nobel --- va lam dieu do hai lan, voi hai linh vuc khoa hoc khac nhau.\\
\textit{(Marie became the first woman to win a Nobel Prize --- and did so twice, in two different fields of science.)}

Bronya luon duoc nhac den trong nhung bai phat bieu cua Marie --- vi khong co su ban be trung thanh ay, khong co Marie Curie.\\
\textit{(Bronya was always mentioned in Marie's speeches --- because without that loyal friendship, there would be no Marie Curie.)}

\end{truyen}
\begin{baihoc}
  Ban be that su dau tu vao nhau --- du khong biet bao gio se thu hoi duoc --- va chinh su dau tu do tao nen vi dai.\\
  \textit{(True friends invest in each other --- not knowing when or whether there will be a return --- and that investment creates greatness.)}
\end{baihoc}

\section{Luu Binh Va Duong Le --- Hung Co Biet Chia Sot Kho}
\begin{truyen}{Luu Binh Va Duong Le --- Hung Co Biet Chia Sot Kho}{Truyen Co Viet Nam --- TK 17}
\chuhoa{L}{}L u Binh va Duong Le la ban than tu thuon, cung nhau an hoc ngheo kho trong nhieu nam.\\
\textit{(Luu Binh and Duong Le were close friends since youth, studying together through years of poverty.)}

Duong Le thi dau len lam quan; Luu Binh thi mai that bai, phai song ban ray lam muon.\\
\textit{(Duong Le passed the exam and became an official; Luu Binh kept failing and had to live hand to mouth.)}

Luu Binh thi den nha ban xin giup do, nghi rang ban se mo long; nhung Duong Le tiep don lanh nhat, co tinh to ra khinh khi.\\
\textit{(Luu Binh traveled to his friend's home hoping for help, expecting a warm welcome; but Duong Le received him coldly, deliberately showing contempt.)}

Bi ton thuong va tuc gian, Luu Binh quyet tam lai sach, cuoi cung thi dau dung truoc mat thien ha.\\
\textit{(Hurt and furious, Luu Binh resolved to study hard, and finally passed the exam with honors before everyone.)}

Chi luc nay Duong Le moi tiet lo: ong da gia vo lanh lung vi biet chi su tuyet vong moi khai hoa su quyet tam cua ban.\\
\textit{(Only then did Duong Le reveal: he had pretended coldness knowing that only the pain of rejection could unlock his friend's resolve.)}

\end{truyen}
\begin{baihoc}
  Co nhung ban be hy sinh su de chiu cua ban than de thuc day ban tien --- chinh ho moi la ban tre long nhat.\\
  \textit{(Some friends sacrifice their own comfort to push you forward --- they are the truest friends of all.)}
\end{baihoc}

\section{Alexander Dai De Va Hephaestion}
\begin{truyen}{Alexander Dai De Va Hephaestion}{Alexander --- Hy Lap TK 4 TCN}
\chuhoa{A}{}A lexander Dai De --- nguoi chinh phuc ba chau luc khi moi 32 tuoi --- co mot nguoi ban thiet than suot cuoc doi la Hephaestion.\\
\textit{(Alexander the Great --- who conquered three continents by age 32 --- had one lifelong friend: Hephaestion.)}

Ho lon len cung nhau, hoc cung Aristotle, chien dau sai nhau trong moi tran ma.\\
\textit{(They grew up together, studied under Aristotle, fought side by side in every battle.)}

Khi mot phai vien Ba Tu nham lam Ho la Alexander va cui dau, Alexander noi: 'Khong sao --- Ho cung la Alexander.'\\
\textit{(When a Persian envoy mistakenly bowed to Hephaestion thinking him Alexander, Alexander said: 'No matter --- he is also Alexander.')}

Khi Hephaestion chet dot ngot o Ecbatana, Alexander khoc long rong --- roi ra lenh quoc tang trong ba ngay.\\
\textit{(When Hephaestion died suddenly at Ecbatana, Alexander wept publicly --- then ordered three days of state mourning.)}

Nguoi ta noi Alexander sau do chang con vui the nua --- chien thang mat vi sac, boi vi thieu nguoi de chia se no.\\
\textit{(People say Alexander was never the same again --- victories lost their flavor, because there was no one left to share them with.)}

\end{truyen}
\begin{baihoc}
  Chua du suc manh de co ban be tri ky --- khi co roi, dung de mat --- vi chien thang khong co ai chia se chi la su co don co doi.\\
  \textit{(Having the strength to find a true friend is one thing --- once found, don't lose them --- for victory shared with no one is only gilded loneliness.)}
\end{baihoc}

\section{Abraham Lincoln Va Joshua Speed}
\begin{truyen}{Abraham Lincoln Va Joshua Speed}{Lincoln --- Hoa Ky 1837}
\chuhoa{K}{}K hi Abraham Lincoln den Springfield hoc luat vao nam 1837, ong gan nhu khong co gi --- khong tien, khong noi o.\\
\textit{(When Abraham Lincoln arrived in Springfield to study law in 1837, he had almost nothing --- no money, no place to stay.)}

Joshua Speed, len noi, nghe chuyen, roi don gian noi: 'Toi co mot giuong kep --- neu ong muon, moi o cung.'\\
\textit{(Joshua Speed, the storekeeper, listened, then simply said: 'I have a large double bed --- if you'd like, you're welcome to share it.')}

Lincoln leo len gac, dat tui sach xuong va noi: '--- Toi o day roi.'\\
\textit{(Lincoln climbed upstairs, set down his bag, and said: 'Well --- I live here now.')}

Ho o cung bay nam; Speed la nguoi Lincoln viet thu khi bi con troi u am nhat; Speed bay den ben khi co can.\\
\textit{(They lived together for seven years; Speed was the one Lincoln wrote to in his darkest moods; Speed came whenever needed.)}

Lincoln ke: 'Speed la nguoi duy nhat toi ke ra ma lo au cua minh va nghe ro loi khuyen tra lai.'\\
\textit{(Lincoln said: 'Speed was the only one to whom I told my anxieties clearly and whose advice I could actually hear.')}

\end{truyen}
\begin{baihoc}
  Ban be that su khong bo di luc ban o tay trang --- ho den --- va o lai; chinh dieu do tao nen su khac biet.\\
  \textit{(True friends don't leave when you're empty-handed --- they arrive --- and stay; that is precisely what makes all the difference.)}
\end{baihoc}

\section{Nguyen Du Va Tinh Tri Am Qua He Thong Chuyen}
\begin{truyen}{Nguyen Du Va Tinh Tri Am Qua He Thong Chuyen}{Nguyen Du --- Viet Nam TK 18-19}
\chuhoa{N}{}N guyen Du song trong thoi ky chuyen tiep day bao dong giua nha Le suy tan va trieu Nguyen moi noi.\\
\textit{(Nguyen Du lived in a tumultuous transitional era between the declining Le dynasty and the rising Nguyen.)}

Ong viet Truyen Kieu --- khong chi la tac pham van hoc ma con la tho dan to am thanh cua trai tim co don.\\
\textit{(He wrote the Tale of Kieu --- not just a literary work but a love poem to the sound of a lonely heart.)}

'Tran the bach nien cuong ban le, De thuong nhung luong nguai biet nhau' --- cau tho noi len noi khat khao tri am.\\
\textit{('In a hundred years of this world, how many people do we truly meet who understand us?' --- this line voices the craving for a soul-friend.)}

Nguyen Du tin rang su hieu nhau that su qua chuyen --- qua van chuong --- con lau ben hon tat ca cac rang buoc xa hoi.\\
\textit{(Nguyen Du believed that true understanding through story --- through literature --- outlasts all social bonds.)}

Truyen Kieu ngay nay duoc doc va kien giai boi hang trieu nguoi --- la minh chung rang tri am co the vuot thoi gian.\\
\textit{(The Tale of Kieu is today read and interpreted by millions --- proof that a soul connection can transcend time.)}

\end{truyen}
\begin{baihoc}
  Van chuong lon la mot loai ban be --- no hieu ban ma khong can ban giai thich; no o lai khi tat ca nguoi khac di.\\
  \textit{(Great literature is a kind of friendship --- it understands you without needing explanation; it stays when everyone else has gone.)}
\end{baihoc}

\section{David Va Jonatan --- Tinh Ban Vuot Quyen Loi}
\begin{truyen}{David Va Jonatan --- Tinh Ban Vuot Quyen Loi}{Kinh Thanh --- Israel TK 10 TCN}
\chuhoa{J}{}J onatan --- con trai vua Sau-lol --- le ra se la nguoi ke vi --- nhung ong yeu quy David hon ca ngai bau.\\
\textit{(Jonathan --- son of King Saul --- was the crown prince --- but he cherished David more than the throne itself.)}

Khi vua cha dat ke hoach giet David, Jonatan tiet lo bi mat va giup David tron thoat.\\
\textit{(When his royal father planned to kill David, Jonathan revealed the secret and helped David escape.)}

Ong tra ve chiu de bat dua cua cha --- trong khi biet rang viec lam cua minh co the bi coi la phan boi.\\
\textit{(He returned to endure his father's rage --- knowing his action could be considered treason.)}

Ong noi voi David: 'Di di --- van menh may se lon lao; chic ta khong the o ben may. Nhung ta se luon la ban may.'\\
\textit{(He said to David: 'Go --- your destiny will be great; I cannot be by your side. But I will always be your friend.')}

Tien su ke: 'Tinh ban cua Jonatan voi David vuot qua tinh yeu cua phu nu' --- vi no khong mang dieu kin nao ngoai su hi sinh.\\
\textit{(Scripture records: 'Jonathan's love for David surpassed the love of women' --- because it carried no condition except self-sacrifice.)}

\end{truyen}
\begin{baihoc}
  Ban be that su khong ganh ghet thanh cong cua ban --- ho vui mung cung ban, dan du no lam mo nhat chan troi cua chinh ho.\\
  \textit{(True friends don't envy your success --- they rejoice with you, even if it dims their own horizon.)}
\end{baihoc}

\section{Goethe Va Schiller --- Doi Ban Biet Nhau Qua Trang Giay}
\begin{truyen}{Goethe Va Schiller --- Doi Ban Biet Nhau Qua Trang Giay}{Goethe & Schiller --- Duc 1794}
\chuhoa{G}{}G oethe va Schiller gap nhau o Weimar nam 1794 --- hai nha tho, hai nha tu tuong, luc dau cung khong qua thich nhau.\\
\textit{(Goethe and Schiller met in Weimar in 1794 --- two poets, two thinkers who at first didn't particularly like each other.)}

Nhung qua trao doi thu tu --- hang tram buc thu trong muoi nam --- ho xay dung nen su hieu biet sau sac.\\
\textit{(But through an exchange of letters --- hundreds of letters over ten years --- they built a profound mutual understanding.)}

Goethe viet ve khoa hoc va triet hoc; Schiller viet ve my hoc va lich su; ca hai cham chua nhau.\\
\textit{(Goethe wrote about science and philosophy; Schiller wrote about aesthetics and history; each critiqued and enriched the other.)}

Nhung tac pham lon nhat cua ca hai --- Faust, Don Carlos, Wilhelm Tell --- duoc viet trong giai doan ho quen biet nhau.\\
\textit{(The greatest works of both --- Faust, Don Carlos, William Tell --- were written during the years they knew each other.)}

Khi Schiller mat som o tuoi 45, Goethe noi: 'Toi mat mot nua ban than minh.'\\
\textit{(When Schiller died young at 45, Goethe said: 'I have lost half of myself.')}

\end{truyen}
\begin{baihoc}
  Ban tri ky khong chi la nguoi chia vui se buon --- ho la nguoi giup ban tro nen nguoi tot hon qua su phan bien va khich le.\\
  \textit{(A soul-friend is not just someone who shares joys and sorrows --- they help you become better through critique and encouragement.)}
\end{baihoc}

\section{Franklin Roosevelt Va Louis Howe}
\begin{truyen}{Franklin Roosevelt Va Louis Howe}{FDR & Louis Howe --- Hoa Ky 1910}
\chuhoa{K}{}K hi Franklin Roosevelt moi vao chinh truong, Louis Howe --- phong vien gia, hom hoi --- quyet dinh ong nay se la Tong Thong.\\
\textit{(When Franklin Roosevelt first entered politics, Louis Howe --- an aging, sickly reporter --- decided this man would be President.)}

Howe chiu trach nhiem moi viec --- tu viet dien van toi ky hop dong --- moi thu de Roosevelt toa sang.\\
\textit{(Howe took charge of everything --- from writing speeches to signing contracts --- everything so Roosevelt could shine.)}

Khi Roosevelt bi liet hai chan vi benh bai liet, Howe la nguoi duy nhat noi: 'Ong van co the thanh Tong Thong.'\\
\textit{(When Roosevelt was paralyzed by polio, Howe was the only one who said: 'You can still be President.')}

Howe khong bao gio tu nhan minh la nguoi quan trong --- ong tu goi minh la 'nguoi hau can' cua Roosevelt.\\
\textit{(Howe never claimed importance for himself --- he called himself Roosevelt's 'errand boy.')}

FDR tro thanh Tong Thong duy nhat trong lich su My giu 4 nhiem ky --- Louis Howe la nguoi o sau hau truong giup dieu do.\\
\textit{(FDR became the only President in American history to serve 4 terms --- Louis Howe was the man behind the curtain who made that possible.)}

\end{truyen}
\begin{baihoc}
  Ban be that su doi khi hi sinh ca su nghiep ca nhan de giup ban thanh cong --- do khong phai su hy sinh ma la tinh yeu.\\
  \textit{(True friends sometimes sacrifice their own career to help you succeed --- that is not sacrifice, that is love.)}
\end{baihoc}
